\documentclass[11pt,a4paper]{article}

% --- Paquetes ---
\usepackage[utf8]{inputenc}
\usepackage[T1]{fontenc}
\usepackage[spanish]{babel}
\usepackage{amsmath, amssymb, amsthm}
\usepackage{geometry}
\usepackage{tikz}
\usepackage{pgfplots}
\usepackage{enumitem}
\usepackage{hyperref}
\usepackage{xcolor}
\usepackage{tcolorbox}

\geometry{margin=2.5cm}
\pgfplotsset{compat=1.18}

% --- Definiciones ---
\newcommand{\barvar}[1]{\overline{#1}}

\newtcolorbox{definicionbox}[1][]{colback=yellow!10, colframe=red!70!black, title=#1, fonttitle=\bfseries}
\newtcolorbox{resultadobox}[1][]{colback=yellow!10, colframe=red!80!black, title=#1, fonttitle=\bfseries}
\newtcolorbox{explicacionbox}[1][]{colback=red!5, colframe=red!60!black, title=#1, fonttitle=\bfseries}
\newtcolorbox{notabox}[1][]{colback=pink!10, colframe=red!50!black, title=#1, fonttitle=\bfseries}

% --- Documento ---
\title{\textbf{Clase 10: La Regla de Oro en el Modelo de Solow} \\ \large Macroeconomía II}
\author{Universidad Católica del Norte}
\date{}

\begin{document}
\maketitle
\tableofcontents
\newpage

% ============================================================
\section{La Función de Producción Cobb-Douglas}
% ============================================================

Trabajaremos con la función de producción Cobb-Douglas:

\begin{resultadobox}[Función de Producción Cobb-Douglas]
\begin{equation}
    Y = A \, K^{\alpha} L^{\beta}, \qquad \alpha + \beta = 1
\end{equation}
\end{resultadobox}

\begin{notabox}[Supuesto]
Hasta el momento hemos asumido que $A = 1$.
\end{notabox}

% ============================================================
\section{(a) Comprobación de los supuestos de la función de producción}
% ============================================================

\subsection{Rendimientos marginales decrecientes}

\subsubsection{Respecto al capital ($K$)}

Buscamos verificar que $\dfrac{\partial Y}{\partial K} > 0$ y $\dfrac{\partial^2 Y}{\partial K^2} < 0$.

\begin{align}
    \frac{\partial Y}{\partial K} &= \alpha A K^{\alpha-1} L^{\beta} = \frac{\alpha A L^{\beta}}{K^{1-\alpha}} > 0 \\[6pt]
    \frac{\partial^2 Y}{\partial K^2} &= (\alpha-1)\,\alpha A K^{\alpha-2} L^{\beta} = \frac{(\alpha-1)\,\alpha A L^{\beta}}{K^{2-\alpha}} < 0
\end{align}

La segunda derivada es negativa porque $(\alpha - 1) < 0$ (ya que $0 < \alpha < 1$), mientras que el resto de los términos son positivos. El producto marginal del capital es positivo y decreciente.

\subsubsection{Respecto al trabajo ($L$)}

Buscamos verificar que $\dfrac{\partial Y}{\partial L} > 0$ y $\dfrac{\partial^2 Y}{\partial L^2} < 0$.

\begin{align}
    \frac{\partial Y}{\partial L} &= \beta A K^{\alpha} L^{\beta-1} > 0 \\[6pt]
    \frac{\partial^2 Y}{\partial L^2} &= (\beta-1)\, A K^{\alpha} L^{\beta-2} < 0
\end{align}

Nuevamente, la segunda derivada es negativa porque $(\beta - 1) < 0$. El producto marginal del trabajo es positivo y decreciente.

\subsection{Rendimientos constantes a escala}

Una función presenta rendimientos constantes a escala si al multiplicar todos los factores por $\lambda > 0$, la producción se multiplica por el mismo factor. Usamos $\lambda = 2$ como ejemplo:

\begin{align}
    F(2K,\, 2L) &= A (2K)^{\alpha}(2L)^{\beta} \notag \\
    &= A \, 2^{\alpha+\beta} K^{\alpha} L^{\beta} \notag \\
    &= 2 \cdot \left( A K^{\alpha} L^{\beta}\right) \notag \\
    &= 2 \cdot F(K,L) = 2 \cdot Y
\end{align}

Donde se usó que $\alpha + \beta = 1$, por lo que $2^{\alpha+\beta} = 2^{1} = 2$. La función exhibe rendimientos constantes a escala.

\subsection{Condiciones de Inada}

\subsubsection{El producto marginal tiende a infinito cuando el factor tiende a cero}

\[
\lim_{K \to 0} \frac{\partial Y}{\partial K} = \lim_{L \to 0} \frac{\partial Y}{\partial L} = \infty
\]

\begin{align}
    \lim_{K \to 0} \frac{\alpha A L^{\beta}}{K^{1-\alpha}} &= \alpha A \lim_{K \to 0} \frac{L^{\beta}}{K^{1-\alpha}} = \alpha A \cdot \infty = \infty \\[6pt]
    \lim_{L \to 0} \frac{\beta A K^{\alpha}}{L^{1-\beta}} &= \beta A \lim_{L \to 0} \frac{K^{\alpha}}{L^{1-\beta}} = \beta A \cdot \infty = \infty
\end{align}

\subsubsection{El producto marginal tiende a cero cuando el factor tiende a infinito}

\[
\lim_{K \to \infty} \frac{\partial Y}{\partial K} = \lim_{L \to \infty} \frac{\partial Y}{\partial L} = 0
\]

\begin{align}
    \lim_{K \to \infty} \frac{\alpha A L^{\beta}}{K^{1-\alpha}} &= \alpha A \lim_{K \to \infty} \frac{L^{\beta}}{K^{1-\alpha}} = \alpha A \cdot 0 = 0 \\[6pt]
    \lim_{L \to \infty} \frac{\beta A K^{\alpha}}{L^{1-\beta}} &= \beta A \lim_{L \to \infty} \frac{K^{\alpha}}{L^{1-\beta}} = \beta A \cdot 0 = 0
\end{align}

\subsection{Esencialidad de los factores}

Ningún factor por sí solo puede producir: si un factor es nulo, la producción es nula.

\begin{align}
    F(0,\, L) &= A \cdot (0)^{\alpha} L^{\beta} = 0 \\[6pt]
    F(K,\, 0) &= A \cdot K^{\alpha} (0)^{\beta} = 0
\end{align}

% ============================================================
\section{(b) Función de producción per cápita}
% ============================================================

Partimos de la función de producción y la dividimos por $L$ (es decir, multiplicamos por $\tfrac{1}{L}$) para expresarla en términos por trabajador:

\begin{align}
    y \;=\; \frac{Y}{L} &= \frac{A K^{\alpha} L^{\beta}}{L} = \frac{A K^{\alpha} L^{\beta}}{L^{\alpha} L^{\beta}} \notag \\[6pt]
    &= A \, \frac{K^{\alpha}}{L^{\alpha}} = A \left(\frac{K}{L}\right)^{\alpha}
\end{align}

Donde se usó que $L = L^{\alpha + \beta} = L^{\alpha} L^{\beta}$. Definiendo el capital per cápita $k \equiv \tfrac{K}{L}$:

\begin{resultadobox}[Función de producción per cápita]
\begin{equation}
    y = A \, k^{\alpha} \;\equiv\; f(k)
\end{equation}
\end{resultadobox}

% ============================================================
\section{(c) Nivel de capital de Estado Estacionario}
% ============================================================

\begin{definicionbox}[Condición de Equilibrio]
En el Estado Estacionario el capital per cápita no cambia:
\[
\Delta k = 0
\]
\end{definicionbox}

Partimos de la \textbf{ecuación de acumulación de capital}:

\begin{align}
    \Delta k &= s \cdot f(k) - \delta k = 0 \notag \\[4pt]
    s \cdot f(k) &= \delta k
\end{align}

Reemplazando la función de producción per cápita $f(k) = A k^{\alpha}$:

\begin{align}
    s \cdot A \, k^{\alpha} &= \delta k \notag \\[4pt]
    s \cdot A &= \delta \, \frac{k}{k^{\alpha}} = \delta \, k^{1-\alpha} \notag \\[4pt]
    k^{1-\alpha} &= \frac{s \cdot A}{\delta}
\end{align}

\begin{resultadobox}[Capital de Estado Estacionario]
\begin{equation}
    k^{*} = \left(\frac{s \cdot A}{\delta}\right)^{\frac{1}{1-\alpha}}
\end{equation}
\end{resultadobox}

\subsection{Estática comparativa}

\begin{explicacionbox}[Efectos sobre el Estado Estacionario]
\begin{itemize}
    \item $+$ Ahorro $(s)$ $\;\Rightarrow\; \uparrow k \;\Rightarrow\; \uparrow y$
    \item $+$ Tecnología $(A)$ $\;\Rightarrow\; \uparrow k \;\Rightarrow\; \uparrow y$
    \item $+$ Depreciación $(\delta)$ $\;\Rightarrow\; \downarrow k \;\Rightarrow\; \downarrow y$
\end{itemize}
\end{explicacionbox}

\subsection{Gráfico del Estado Estacionario}

El Estado Estacionario se alcanza donde la inversión $i = s \cdot A k^{\alpha}$ iguala a la depreciación $\delta k$.

\begin{center}
\begin{tikzpicture}
    \begin{axis}[
        axis lines=middle,
        xlabel={$k$}, ylabel={$y$},
        xmin=0, xmax=10, ymin=0, ymax=6,
        xtick=\empty, ytick=\empty,
        every axis x label/.style={at={(ticklabel* cs:1)}, anchor=west},
        every axis y label/.style={at={(ticklabel* cs:1)}, anchor=south},
        width=12cm, height=8cm,
        samples=200,
    ]
    % Función de producción y = A k^alpha
    \addplot[domain=0:10, thick, black] {1.8*x^0.5};
    \node[right, black] at (axis cs:9.5,5.4) {$y = A k^{\alpha}$};

    % Depreciación delta*k
    \addplot[domain=0:9, thick, orange] {0.62*x};
    \node[right, orange] at (axis cs:8.6,5.5) {$\delta k$};

    % Inversión i = s*A*k^alpha
    \addplot[domain=0:10, thick, violet] {0.9*x^0.5};
    \node[right, violet] at (axis cs:9.5,2.95) {$i = s \cdot A k^{\alpha}$};

    % Capital de estado estacionario (corte entre i y delta*k): 0.9*sqrt(x)=0.62*x -> x = (0.9/0.62)^2
    \addplot[only marks, mark=*, mark size=2pt, black] coordinates {(2.107,1.307)};
    \addplot[dashed, gray] coordinates {(2.107,0) (2.107,1.307)};
    \node[below] at (axis cs:2.107,0) {$k^{*}=\left(\frac{s \cdot A}{\delta}\right)^{\frac{1}{1-\alpha}}$};

    \end{axis}
\end{tikzpicture}
\end{center}

% ============================================================
\section{(d) El capital de la Regla de Oro}
% ============================================================

\begin{definicionbox}[Regla de Oro]
El capital de la Regla de Oro es el nivel de capital per cápita que \textbf{maximiza el consumo} per cápita en el Estado Estacionario.
\end{definicionbox}

\subsection{Problema del planificador}

El planificador elige el capital $k^{*}$ que maximiza el consumo de Estado Estacionario:

\[
\max_{k^{*}} \; c^{*}(k^{*}) = f(k^{*}) - \delta k^{*}
\]

\begin{notabox}[¿Por qué $f(k^{*}) - \delta k^{*}$?]
Porque en el Estado Estacionario la inversión iguala a la depreciación:
\[
\delta k^{*} = s \cdot f(k^{*})
\]
de modo que el consumo es lo que queda del producto luego de reponer el capital depreciado.
\end{notabox}

Reemplazando $f(k^{*}) = A k^{*\alpha}$:

\[
\max_{k^{*}} \; c^{*}(k^{*}) = A k^{*\alpha} - \delta k^{*}
\]

\subsection{Condición de primer orden (FOC)}

\begin{align}
    \frac{d c^{*}}{d k^{*}} &= \alpha A k^{*\,\alpha-1} - \delta = 0 \notag \\[6pt]
    \underbrace{\alpha A k^{*\,\alpha-1}}_{\text{PMK}} &= \delta
\end{align}

\begin{resultadobox}[Condición de la Regla de Oro]
\begin{equation}
    \text{PMK} = \delta
\end{equation}
El consumo se maximiza cuando el producto marginal del capital iguala a la tasa de depreciación.
\end{resultadobox}

\subsection{Despeje del capital de la Regla de Oro}

\begin{align}
    \alpha A k^{*\,\alpha-1} &= \delta \notag \\[4pt]
    k^{*\,\alpha-1} &= \frac{\delta}{\alpha A} \notag \\[4pt]
    k^{*\,1-\alpha} &= \frac{\alpha A}{\delta}
\end{align}

\begin{resultadobox}[Capital de la Regla de Oro]
\begin{equation}
    k^{*}_{\text{GOLD}} = \left(\frac{\alpha \cdot A}{\delta}\right)^{\frac{1}{1-\alpha}}
\end{equation}
\end{resultadobox}

\end{document}
